\section{Conclusion}
  \label{sec:conclusion}

  We have constructed a complete explicit example of spectral
  admissibility filtering on the $Q_8$ Cayley graph.

  The analysis establishes three structural results.

  First, the bounded-flux constraint
  $|\partial_t \chi| \le c_\chi$
  induces a hierarchy of maximal admissible effective amplitudes
  \[
    A_n^{\max} = \frac{c_\chi}{\sqrt{\lambda_n}},
  \]
  favouring the non-abelian quaternionic sector over the
  nontrivial abelian sectors.

  Second, the hierarchy remains stable under the weakly
  non-linear Dirac–Born–Infeld correction.
  The leading amplitude-dependent frequency shift depends
  only on the Laplacian eigenvalue $\lambda_n$.

  Third, the admissibility envelope coincides with the
  exact saturation boundary of the DBI dynamics.
  The hierarchy therefore persists beyond the perturbative
  regime.

  The origin of this behaviour can be traced to the algebraic
  structure of the Cayley graph.
  The quaternionic sector satisfies $\mu_{\rho_4}=0$,
  placing it at the centre of the adjacency spectrum and
  producing the smallest non-trivial Laplacian eigenvalue.

  This explicit computation provides a minimal demonstration
  that bounded relational capacity can favour non-abelian
  spectral sectors purely through admissibility filtering.

  More generally, the mechanism suggests that non-abelian
  symmetry sectors may correspond to dynamically resilient
  regions of the relational spectrum under bounded-flux
  constraints.

  The $Q_8$ graph represents the smallest non-trivial
  example of this mechanism.
  Larger Cayley graphs approximating $S^3$, such as those
  associated with the binary icosahedral group $2I$,
  provide natural extensions for investigating whether
  the same admissibility hierarchy persists in larger
  relational substrates.
